\section{Problem Definition}
\label{sec:problem}
This section defines the attribute synthesis problem and the correctness
requirements for its solutions. Given a rewrite's LHS metadata, the goal is
to construct RHS attributes that make the replacement legal and preserve
its result under the declared conditions. We also specify the mathematical
fragment for which we establish correctness and completeness.

\subsection{Tensor Expressions and Rewrites}
Let $\mathcal X=\{X_1,\ldots,X_n\}$ be the input tensor names. Each $X_i$ has
symbolic shape $s_i$; a match supplies its concrete shape. Let $\mathcal O$
be the operations. An operation $o\in\mathcal O$ has $k$ inputs and an
attribute descriptor $a$. Tensor expressions have the form
\begin{equation}
 T ::= X_i \mid o_a(T_1,\ldots,T_k).
\end{equation}
A rewrite is $L\Rightarrow R$. Its LHS attributes are known, and its RHS
attributes are to be constructed. We use the following notation throughout:
\begin{equation}
\begin{aligned}
 s&=(s_1,\ldots,s_n) &&\text{input shapes},\\
 \alpha&=\text{LHS attributes},\\
 p&=(s,\alpha) &&\text{source metadata},\\
 \beta&=\text{RHS attributes}.
\end{aligned}
\label{eq:metadata}
\end{equation}
Metadata $p$ is concrete at a match and makes every LHS operation legal.
Writing $X=(X_1,\ldots,X_n)$ for the input values, the two computations are
$L_p(X)$ and $R_{p,\beta}(X)$. An attribute program reads $p$, never the
entries of $X$.

\subsection{Valid Attribute Assignments}
\label{sec:valid}
An assignment must make every RHS operation legal and preserve both ordered
output shape and values:
\begin{equation}
 \operatorname{Valid}(p,\beta)\iff
 \operatorname{Legal}_R(p,\beta)
 \land\operatorname{Shape}_L(p)=\operatorname{Shape}_R(p,\beta)
 \land\forall X,\; L_p(X)=R_{p,\beta}(X).
\label{eq:valid}
\end{equation}
Here $X$ ranges over all values with the shapes in $p$.
$\operatorname{Legal}_R$ includes every intermediate operation, and a shape
is an ordered vector. The same $\beta$ must work for every compatible $X$:
the quantifier order is $\exists\beta\,\forall X$, not a separate attribute
choice for each tensor value.

\subsection{What Must Synthesis Produce?}
\label{sec:find}
Our goal is to synthesize two reusable programs, \code{check} and
\code{build}, from a symbolic rewrite. Given source metadata $p$,
\code{check} decides whether valid RHS attributes exist, and \code{build}
constructs them whenever \code{check} succeeds. The programs must satisfy:
\begin{align}
 \operatorname{check}(p)&\iff\exists\beta:\operatorname{Valid}(p,\beta),
 \label{eq:check-goal}\\
 \operatorname{check}(p)&\Rightarrow
 \operatorname{Valid}(p,\operatorname{build}(p)).
 \label{eq:build-goal}
\end{align}
Both programs depend only on source metadata and must satisfy these
requirements for every legal instance in the defined fragment, across
varying ranks and shapes.

\subsection{Scope and Assumptions of the Completeness Theorem}
\label{sec:formal-scope}
Sections~\ref{sec:axis-model} and~\ref{sec:mathematical-synthesis} use the following
assumptions. TEPL also handles operations and rules outside this model;
Section~\ref{sec:broader-operators} explains that separate path.

\paragraph{Operation contracts.}
The proof uses exactly the operation contracts in
Table~\ref{tab:formal-operations}. Empty attribute lists are permitted.
\begin{table}[!htbp]
\caption{Operations in the mathematical model.}
\label{tab:formal-operations}
\begin{tabularx}{\linewidth}{@{}p{0.17\linewidth}YY@{}}\toprule
Operation & Attributes and legality & Computation and output\\\midrule
Permutation & A full permutation $\pi$ of input positions.
& Output axis $j$ selects input axis $\pi_j$.\\
Unary\newline summation & Distinct, in-range positions $d$.
& Sum over $d$, remove those axes, and preserve the remaining axis order.\\
Contraction & Ordered batch lists $(b_\ell,b_r)$ and contraction lists $(k_\ell,k_r)$.
For each operand, all listed positions are distinct and in range.
Corresponding left/right lists have equal lengths and pair axes of equal extent.
& Multiply at paired coordinates and sum over contraction axes. Output order:
batch-list order, then left-free axes, then right-free axes; free axes retain
operand order.\\
Addition & No attributes. Operands have identical ordered shapes; no broadcasting.
& Add corresponding elements and preserve the operand shape and axis order.\\\bottomrule
\end{tabularx}
\end{table}

\paragraph{Inputs and arithmetic.}
Inputs have arbitrary finite ranks, including zero, and axis extents
of at least two. Input values vary independently over the reals,
and arithmetic is exact.

\paragraph{Expression structure.}
A \emph{read-once component} is a contraction/permutation/unary-summation
tree that uses every tensor in a nonempty input set exactly once. The LHS
and RHS of a component use the same input set.

For addition, expand the LHS using linearity and bilinearity into
$L=\sum_t L_t$. The actual RHS must already be an explicit sum of matching
components, $R=\sum_t R_t$. Each $L_t,R_t$ is read-once and uses the same
input set $U_t$. Different terms have different input sets, although these
sets may overlap.

\paragraph{RHS attributes.}
Each RHS operation has its own attributes for the compiler to infer,
subject only to Table~\ref{tab:formal-operations}. The theorem excludes
attributes fixed in advance or required to be shared between RHS operations.
Expanding expressions must preserve captured LHS attributes and their
sharing, without splitting a shared RHS attribute into separate choices.

The completeness theorem applies only to this fragment.
Section~\ref{sec:limitations} discusses cases outside its scope.

\subsection{The Contract Exposed by TEPL}
\label{sec:contract}
Whenever a legal computation matches the LHS and satisfies \code{where},
the rule's attribute program must successfully compute RHS attributes that
produce a legal, equivalent replacement. This applies whether the program
is written explicitly or inferred. Writing $C(p)$ for the condition and $f$
for the attribute program, TEPL seeks to establish
\begin{equation}
 \forall p.\quad \operatorname{Legal}_L(p)\land C(p)
 \Longrightarrow p\in\operatorname{dom}(f)
 \land\operatorname{Valid}(p,f(p)).
\label{eq:contract}
\end{equation}
